\expandafter\gdef\csname SNside1\endcsname{\node[anchor=north west,inner sep=0pt,text width=48mm,font=\SideFont,align=left] at ([xshift=138mm,yshift=-516.15pt]current page.north west) {\textcolor{SNink}{\textbf{最小反例 · 最高位1并不总是负数}}\par\vspace{1.5mm}最小反例：同一8位串FFH按无符号解释是255，按补码解释是\ensuremath{-}1。编码和位宽必须与位串一同给出。直接把最高位称作负号，会把无符号扩展、浮点阶码和地址都误读成有符号数。\par\smallskip\begin{tikzpicture}[x=1mm,y=1mm,>=Stealth,node distance=5mm]
\node[draw=SNink!50,fill=SNblue,rounded corners=1mm,text width=34mm,align=center,inner sep=1.2mm,font=\fontsize{7.5}{9}\selectfont] (m0) at (0,0) {8位FFH};
\node[draw=SNink!50,fill=SNblue,rounded corners=1mm,text width=34mm,align=center,inner sep=1.2mm,font=\fontsize{7.5}{9}\selectfont] (m1) at (0,-11) {无符号255};
\node[draw=SNink!50,fill=SNblue,rounded corners=1mm,text width=34mm,align=center,inner sep=1.2mm,font=\fontsize{7.5}{9}\selectfont] (m2) at (0,-22) {补码−1};
\draw[->,SNink] (m0) -- (m1);
\draw[->,SNink] (m0) edge[bend left=40] (m2);
\end{tikzpicture}
\par{\fontsize{7.5}{9}\selectfont 同一位串依编码得到不同真值。}};}
\expandafter\gdef\csname SNside2\endcsname{\node[anchor=north west,inner sep=0pt,text width=48mm,font=\SideFont,align=left] at ([xshift=138mm,yshift=-397.86pt]current page.north west) {\textcolor{SNink}{\textbf{不变量 · 范围来自位权}}\par\vspace{1.5mm}补码最高位权为负，其余位权为正；最小值只有最高位1，最大值只有最高位0。每次推范围先写位权，就能同时解释零唯一、负数多一个，以及最小负数同宽取负不可表示。};}
\expandafter\gdef\csname SNside3\endcsname{\node[anchor=north west,inner sep=0pt,text width=48mm,font=\SideFont,align=left] at ([xshift=138mm,yshift=-95.00pt]current page.north west) {\textcolor{SNink}{\textbf{错项排除 · 扩展看源数解释}}\par\vspace{1.5mm}FFFAH若是unsigned short则为65530，扩宽补0；若是signed short补码则为\ensuremath{-}6，扩宽复制符号。目标名字中出现unsigned不代表先补0：数值转换还需按目标模数处理。};}
\expandafter\gdef\csname SNside4\endcsname{\node[anchor=north west,inner sep=0pt,text width=48mm,font=\SideFont,align=left] at ([xshift=138mm,yshift=-95.00pt]current page.north west) {\textcolor{SNink}{\textbf{最小反例 · 右移不等于向零除法}}\par\vspace{1.5mm}最小反例：8位补码\ensuremath{-}3为FDH，算术右移得到FEH，即\ensuremath{-}2；C整数\ensuremath{-}3/2向零截断为\ensuremath{-}1。两者相等只在额外条件下成立。检查奇数负数能快速暴露把右移无条件替代除法的错误。};}
\expandafter\gdef\csname SNside5\endcsname{\node[anchor=north west,inner sep=0pt,text width=48mm,font=\SideFont,align=left] at ([xshift=138mm,yshift=-95.00pt]current page.north west) {\textcolor{SNink}{\textbf{不变量 · 减法中的一只加一}}\par\vspace{1.5mm}控制Sub=1时第二路逐位取反，最低进位再加1，总体才是A\ensuremath{-}B。若已经送入B的负数补码，又把最低进位设1，就把那一个1重复加入。追踪输入端比只写``取补码''更精确。};}
\expandafter\gdef\csname SNside6\endcsname{\node[anchor=north west,inner sep=0pt,text width=48mm,font=\SideFont,align=left] at ([xshift=138mm,yshift=-95.00pt]current page.north west) {\textcolor{SNink}{\textbf{错项排除 · 进位和有符号溢出分开}}\par\vspace{1.5mm}CF依据无符号加减的进借位，OF依据补码真值范围。同一运算完全可以CF=1而OF=0，或反过来。先确定ISA对减法标志采用借位还是无借位口径，再套公式。};}
\expandafter\gdef\csname SNside7\endcsname{\node[anchor=north west,inner sep=0pt,text width=48mm,font=\SideFont,align=left] at ([xshift=138mm,yshift=-95.00pt]current page.north west) {\textcolor{SNink}{\textbf{课程连接 · 标志接到条件分支}}\par\vspace{1.5mm}指令系统用ZF判断相等，用CF处理无符号大小，用SF与OF的组合处理有符号大小。编译同一个小于号时还要看操作数类型；电路相同，不代表转移条件可以混用。};}
\expandafter\gdef\csname SNside8\endcsname{\node[anchor=north west,inner sep=0pt,text width=48mm,font=\SideFont,align=left] at ([xshift=138mm,yshift=-259.60pt]current page.north west) {\textcolor{SNink}{\textbf{不变量 · 余数必须最终回验}}\par\vspace{1.5mm}无符号除法结束后，同时核N=QD+R及0\ensuremath{\le}R\textless D。恢复法与不恢复法可以出现不同中间余数，但最终条件相同；不恢复法若结束时R为负，常需补做一次加D校正。};}
\expandafter\gdef\csname SNside9\endcsname{\node[anchor=north west,inner sep=0pt,text width=48mm,font=\SideFont,align=left] at ([xshift=138mm,yshift=-144.70pt]current page.north west) {\textcolor{SNink}{\textbf{错项排除 · 字段长度不等于精度}}\par\vspace{1.5mm}binary32尾数字段23位，正规格化数隐含前导1，因此有效精度24位；非规格化数隐含0。把所有编码都当1.F，会把最小正数、零和特殊值混在一起。先判断阶码是否全0或全1。};}
\expandafter\gdef\csname SNside10\endcsname{\node[anchor=north west,inner sep=0pt,text width=48mm,font=\SideFont,align=left] at ([xshift=138mm,yshift=-250.76pt]current page.north west) {\textcolor{SNink}{\textbf{不变量 · 阶码字段先减偏置}}\par\vspace{1.5mm}普通IEEE浮点数的指数是E减偏置，尾数字段另按二进制小数解释。整数移码常见的翻转符号位关系，不能直接套到偏置127的IEEE阶码。编码和解码各走一遍可验证字段。};}
\expandafter\gdef\csname SNside11\endcsname{\node[anchor=north west,inner sep=0pt,text width=48mm,font=\SideFont,align=left] at ([xshift=138mm,yshift=-95.00pt]current page.north west) {\textcolor{SNink}{\textbf{最小反例 · 不溢出仍可不精确}}\par\vspace{1.5mm}最小反例：binary32能精确存2的24次方，却不能精确存它加1。指数远未到上限，有效位已经不足。浮点溢出与舍入损失要分别检查，答案有限并不保证与数学精确值相等。};}
\expandafter\gdef\csname SNside12\endcsname{\node[anchor=north west,inner sep=0pt,text width=48mm,font=\SideFont,align=left] at ([xshift=138mm,yshift=-95.00pt]current page.north west) {\textcolor{SNink}{\textbf{课程连接 · 编程中的循环边界}}\par\vspace{1.5mm}无符号n=0时n\ensuremath{-}1绕到最大值，会改变循环终止；有符号溢出又涉及语言规则。把循环、编码范围、浮点精度三步分开检查，可以说明同一算法在不同类型中为何产生不同结果。};}
\expandafter\gdef\csname SNside13\endcsname{\node[anchor=north west,inner sep=0pt,text width=48mm,font=\SideFont,align=left] at ([xshift=138mm,yshift=-95.00pt]current page.north west) {\textcolor{SNink}{\textbf{不变量 · 字节序和对齐分两步}}\par\vspace{1.5mm}先依成员类型与对齐规则求对象布局，再把每个值拆成字节按大小端放入。字节序不改变结构体填充，也不反转数组元素次序。两步分别核对后，才能回答某地址保存的具体字节。};}
\expandafter\gdef\csname SNside14\endcsname{\node[anchor=north west,inner sep=0pt,text width=48mm,font=\SideFont,align=left] at ([xshift=138mm,yshift=-95.00pt]current page.north west) {\textcolor{SNink}{\textbf{不变量 · 相邻间隔随指数变大}}\par\vspace{1.5mm}正规格化binary32在同一指数区间内间隔是2的e\ensuremath{-}23次方。指数每增1，相邻间隔也翻倍。先算间隔再判某整数是否可表示，比只记``24位精度''更容易处理大整数与舍入中点。};}
\expandafter\gdef\csname SNside15\endcsname{\node[anchor=north west,inner sep=0pt,text width=48mm,font=\SideFont,align=left] at ([xshift=138mm,yshift=-95.00pt]current page.north west) {\textcolor{SNink}{\textbf{课程连接 · 单位衔接存储系统}}\par\vspace{1.5mm}数据宽度、地址单位和块大小会在下一章共同决定偏移位数与容量。按字节编址时4B整数跨4个地址，但同一字节内8个位并不产生8个独立字节地址。每个对数前都先写清以字节还是以字计数。};}
\expandafter\gdef\csname SNside16\endcsname{\node[anchor=north west,inner sep=0pt,text width=48mm,font=\SideFont,align=left] at ([xshift=138mm,yshift=-95.00pt]current page.north west) {\textcolor{SNink}{\textbf{最小反例 · 成员后面仍可能补空隙}}\par\vspace{1.5mm}最小反例：int之后的char只占1B，但接下来的short起点仍需为2的倍数。因此char与short之间可能填充1B。连续声明不代表在内存中毫无间隙；结构体末尾也可能有填充。};}
\AddToShipoutPictureFG{\begin{tikzpicture}[remember picture,overlay]\ifcsname SNside\thepage\endcsname\csname SNside\thepage\endcsname\fi\end{tikzpicture}}